\begin{styletable}

Множество (перечисление)
    & $\set{1, 2, 3}$
    & \verb|\set{1, 2, 3}|
\\

Множество (условие)
    & $\setc{x \in X}{x > 0}$
    & \verb|\setc{x \in X}{x > 0}|
\\

Мощность множества
    & $\card{X}$
    & \verb|\card{X}|
\\

Принадлежность
    & $x \in A$, $A \ni y$
    & \verb|\in|, \verb|\ni|
\\

Не принадлежит
    & $x \not\in A$, $A \not\ni y$
    & \verb|\not\in|, \verb|\not\ni|
\\

Пустое множество
    & $\emptyset$
    & \verb|\emptyset|
\\

Подмножество
    & $A \subset B$, $B \supset A$
    & \verb|\subset|, \verb|\supset|
\\

Не является подмножеством
    & $A \not\subset B$, $B \not\supset A$
    & \verb|\not\subset|, \verb|\not\supset|
\\

Строгое подмножество
    & $A \subsetneq B, B \supsetneq A$
    & \verb|\subsetneq|, \verb|\supsetneq|
\\

Объединение
    & $A \cup B$
    & \verb|\cup|
\\

Пересечение
    & $A \cap B$
    & \verb|\cap|
\\

Объединение системы
    & $\bigcup_{s \in S} A_s$
    & \verb|\bigcup_{s \in S} A_s|
\\

Пересечение системы
    & $\bigcap_{s \in S} A_s$
    & \verb|\bigcap_{s \in S} A_s|
\\

Разность
    & $A \setminus B$
    & \verb|\setminus|
\\

Прямое произведение
    & $A \times B$
    & \verb|\times|
\\

Отрезок
    & $\seg{a}{b}$
    & \verb|\seg{a}{b}|
\\

Интервал
    & $\sego{a}{b}$
    & \verb|\sego{a}{b}|
\\

Полуинтервал (<<справа>>)
    & $\segor{a}{b}$
    & \verb|\segor{a}{b}|
\\

Полуинтервал (<<слева>>)
    & $\segol{a}{b}$
    & \verb|\segol{a}{b}|
\\

Кортеж
    & $\tup{1, 2, 3}$
    & \verb|\tup{1, 2, 3}|
\\

\hline
\end{styletable}

\stylenote{
    Элементы множества внутри фигурных скобок, а также левые и правые концы числовых интервалов отделяются друг от друга запятой.
}

\stylenote{
    Символ $\subset$ обозначает нестрогое включение множеств (то есть допускается равенство).
    Строгое включение записывается с помощью команд \texttt{\textbackslash subsetneq} ($\subsetneq$) и \texttt{\textbackslash supsetneq} ($\supsetneq$).
    Символы $\subseteq$ и $\supseteq$ не используются.
}